\documentclass[../main.tex]{subfiles}
\begin{document}
\begingroup
\pagestyle{empty}
\section*{表記検討メモ(執筆用, 提出対象外)}
本文中の表記揺れを整理するための暫定メモである。

\small
\setlength{\tabcolsep}{5pt}
\renewcommand{\arraystretch}{1.35}
\begin{longtable}{|p{0.28\linewidth}|p{0.28\linewidth}|p{0.34\linewidth}|}
\hline
現在使っている表記 & 迷っている表記 & 意図など \\
\hline
\endfirsthead
\hline
現在使っている表記 & 迷っている表記 & 意図など \\
\hline
\endhead

連続読み出し式データ収集システム & ストリーミングDAQ & 今まで主に「連続読み出し式DAQ」と書いてきたが、物理学会とかだとセッション名は「データ収集」と表記されている。ストリーミングの方がシステムの原理をより正確に表現している気もする。 \\
\hline

オンラインフィルタ & オンラインフィルター & 発音して心地が良いのは後者 \\
\hline

リアルタイム事象識別プロセス & オンラインフィルタ or ハイレベルフィルター & タイトルでは前者だが、論文中では後者で呼ぶ \\
\hline
\end{longtable}

\subsection*{表記チェック項目 (2026-11-12)}

\begin{itemize}
  \item カッコの全角・半角
  \begin{itemize}
    \item 全角カッコは使わない。
  \end{itemize}

  \item 句読点・読点の全角・半角
  \begin{itemize}
    \item 句点・読点は全角を使う。
    \item 文章中のカンマ・ピリオドについては例外的に半角を使うこともある。
  \end{itemize}

  \item 全角英数字
  \begin{itemize}
    \item 全角英数字は使わない。
  \end{itemize}

  \item 用語と長音
  \begin{itemize}
    \item 本文の「ハイレベルフィルタ」と付録10の9行の「ハイレベルフィルター」が混在。
  \end{itemize}

  \item 略語の初出
  \begin{itemize}
    \item DAQ、FEE、TOF、DC、SFT、T0、BFT、LTOF、STLなどについて、意味と必要な初出説明があるかを点検。
    \item 現状は多くが未説明。
  \end{itemize}

  \item 誤記・重複
  \begin{itemize}
    \item 
  \end{itemize}

  \item 見出し・番号・参照
  \begin{itemize}
    \item
  \end{itemize}

  \item 引用・主張の根拠
  \begin{itemize}
    \item
  \end{itemize}

  \item 図表・コードのcaptionとlabel
  \begin{itemize}
    \item
  \end{itemize}

\end{itemize}

\clearpage
\endgroup
\end{document}
